\documentclass{article}

\usepackage{xcolor}
\usepackage{graphicx}
\usepackage{fancyvrb}
\usepackage{listings}
\usepackage[T1]{fontenc}
\usepackage{hyperref}
\usepackage{amsmath}

\definecolor{officegreen}{rgb}{0, 0.5, 0}
\definecolor{navy}{rgb}{0, 0, 0.5}
\definecolor{linecolor}{rgb}{0.5, 0.6875, 0.6875}
\definecolor{outputcolor}{rgb}{0.375, 0.375, 0.375}

\newcommand{\id}[1]{\textcolor{black}{#1}}
\newcommand{\com}[1]{\textcolor{officegreen}{#1}}
\newcommand{\inact}[1]{\textcolor{gray}{#1}}
\newcommand{\kwd}[1]{\textcolor{navy}{#1}}
\newcommand{\num}[1]{\textcolor{officegreen}{#1}}
\newcommand{\ops}[1]{\textcolor{purple}{#1}}
\newcommand{\prep}[1]{\textcolor{purple}{#1}}
\newcommand{\str}[1]{\textcolor{olive}{#1}}
\newcommand{\lines}[1]{\textcolor{linecolor}{#1}}
\newcommand{\fsi}[1]{\textcolor{outputcolor}{#1}}
\newcommand{\omi}[1]{\textcolor{gray}{#1}}

% Overriding color and style of line numbers
\renewcommand{\theFancyVerbLine}{
\lines{\small \arabic{FancyVerbLine}:}}

\lstset{%
  backgroundcolor=\color{gray!15},
  basicstyle=\ttfamily,
  breaklines=true,
  columns=fullflexible
}

\title{{page-title}}
\date{}

\begin{document}

\maketitle


\section*{Example: Mermaid Diagrams}



\href{https://mermaid.js.org/}{Mermaid} is a JavaScript-based diagramming and charting tool that renders Markdown-inspired text definitions into diagrams.


The recommended way to use Mermaid with fsdocs is to write diagrams as plain fenced code blocks, and add a small script that turns those blocks into diagrams when the page loads. The Markdown stays portable: GitHub renders \texttt{```mermaid} fences natively, and your fsdocs site shows the real diagrams. This very page uses the pattern below, so the diagrams you see are fenced code blocks promoted by a \texttt{\_body.html} script.
\subsection*{Setup}



Create or edit a \texttt{\_body.html} file in your \texttt{docs} folder. fsdocs injects it at the end of every page, after the content. The script imports mermaid and promotes fenced mermaid blocks to the \texttt{<div class="mermaid">} elements mermaid looks for:
\begin{lstlisting}
<script type="module">
  import mermaid from 'https://cdn.jsdelivr.net/npm/mermaid@11/dist/mermaid.esm.min.mjs';

  // A ```mermaid fenced block renders natively on GitHub, but reaches fsdocs
  // as a syntax-highlighted code block. Promote those blocks to
  // <div class="mermaid"> elements so one plain fence works in both places.
  for (const code of document.querySelectorAll('code[lang="mermaid"]')) {
    // fsdocs wraps the snippet in <table class="pre"><tr><td><pre><code>...;
    // replace the outermost wrapper so no table scaffolding is left around
    // the diagram.
    const snippet = code.closest('table.pre') ?? code.closest('pre');
    if (!snippet) continue;
    const diagram = document.createElement('div');
    diagram.className = 'mermaid';
    // textContent, not innerHTML: the source arrives HTML-escaped
    // (arrows come through with escaped angle brackets) and mermaid needs
    // the raw arrows back.
    diagram.textContent = code.textContent;
    snippet.replaceWith(diagram);
  }

  mermaid.initialize({ startOnLoad: true });
</script>

\end{lstlisting}
\subsection*{Usage}



Write your diagram in a fenced code block with the \texttt{mermaid} language tag:
\begin{lstlisting}
```mermaid
graph LR
    A[Input docs] --> B[fsdocs build]
    B --> C[HTML output]
    B --> D[API reference]
```

\end{lstlisting}


On this site, the block above is rendered as:
\begin{lstlisting}
graph LR
    A[Input docs] --> B[fsdocs build]
    B --> C[HTML output]
    B --> D[API reference]

\end{lstlisting}
\subsection*{More Examples}



Sequence diagram:
\begin{lstlisting}
sequenceDiagram
    participant User
    participant fsdocs
    participant Browser
    User->>fsdocs: dotnet fsdocs watch
    fsdocs-->>Browser: Serve docs
    User->>fsdocs: Edit .md or .fsx
    fsdocs-->>Browser: Reload page

\end{lstlisting}


Class diagram:
\begin{lstlisting}
classDiagram
    class ApiDocComment {
        +Summary: string
        +Remarks: string option
        +Parameters: ApiDocSection list
    }
    class ApiDocMember {
        +Name: string
        +Comment: ApiDocComment
    }
    ApiDocMember --> ApiDocComment

\end{lstlisting}
\subsection*{Tips}

\begin{itemize}
\item To customise the Mermaid theme, pass options to \texttt{mermaid.initialize()}, for example \texttt{theme: "base"} together with \texttt{themeVariables}.

\item To centre the diagrams, add a rule for the promoted element to your \texttt{docs/content/fsdocs-theme.css}:

\end{itemize}

\begin{lstlisting}
.mermaid {
  margin: 1rem auto;
  & svg {
    display: block;
    margin: 0 auto;
  }
}

\end{lstlisting}
\begin{itemize}
\item You can also write \texttt{<div class="mermaid">} blocks directly in your Markdown. The promotion script only touches fenced code blocks, so both forms can coexist.

\item See the \href{https://mermaid.js.org/intro/}{Mermaid documentation} for the full list of supported diagram types.

\end{itemize}



\end{document}